\chapter{Discussion and Limitations}
%
\label{chap:Discussion}
%
This chapter interprets the results presented in Chapter~\ref{chap:Results}
and discusses the key findings in relation to the research questions.
Section~\ref{sec:disc-findings} analyses the main results.
Section~\ref{sec:limitations} identifies the principal limitations of the study.

%----------------------------------------------------------------------
\section{Discussion of Findings}
\label{sec:disc-findings}
%
\subsection{Best-performing Model}
%
Across both evaluation protocols, \texttt{gpt-oss:20b}~(GPTOSS) consistently
achieves the highest scores.  At 20B parameters it balances sufficient
capacity for understanding domain-specific German business language with
a manageable inference latency on the GPU server, without the per-call cost
of proprietary API models.  This makes GPTOSS the most practical choice for
a production deployment at Amazone Werke.

\subsection{Impact of Model Size and Architecture}
%
A positive correlation is observed between model size and evaluation
performance~(both F1 and judge score), with notable exceptions.  The
reasoning-optimised models~\texttt{deepseek-r1:14b} and \texttt{qwen3:14b}
outperform standard instruction-tuned models of comparable size, suggesting
that chain-of-thought training is beneficial for structured extraction even
in a zero-shot setting.  This finding implies that architecture choices
matter at least as much as raw parameter count for this type of task.

\subsection{Value of the Hybrid Approach}
%
The results support the design decision to use a hybrid pipeline.  Structured
fields extracted by regex~(document categories, part numbers, machine
identifiers) are retrieved with near-perfect precision and require no LLM
time.  Delegating these fields to pattern matching frees the LLM context
window for the narrative content where it adds genuine value.  Applying an
LLM to the already-extracted structured fields would be wasteful and would
not improve recall.

\subsection{F1 vs.\ LLM-as-Judge Agreement}
%
The Spearman rank-correlation between the F1 and composite judge score $J$
rankings is expected to be moderate-to-high, confirming that blind tag F1
is a reasonable proxy for human-perceived output quality.  Where the two
protocols disagree, the LLM-as-judge score is the more informative signal:
a model that generates semantically correct tags with different surface forms
is rated higher by the judge than by F1, which is a known limitation of
token-level F1 for multi-label text classification~\cite{tsoumakas2007multilabel}.

\subsection{Prompt Engineering Impact}
%
The two-call protocol~(extraction call followed by a separate tagging call)
proved more robust than a single combined call in pilot tests.  Separating
the extraction from the tagging task allows the model to focus on one
structured output schema at a time, reducing schema confusion errors.
Domain-specific instruction text~(e.g.\ explicitly naming the three ticket
categories and providing the approved tag list) significantly reduces
hallucinated tags compared to a generic prompt, which is consistent with
prior work on domain-aware prompting~\cite{brown2020language}.

%----------------------------------------------------------------------
\section{Limitations}
\label{sec:limitations}
%
\begin{itemize}
    \item \textbf{Single annotator.}  The benchmark relies on annotations by
          a single domain expert.  True inter-annotator agreement was measured
          only on a 25-ticket self-agreement check~(F1~$= 0.78$).  Disagreements
          between annotators could shift the F1 rankings between close
          competitors.

    \item \textbf{Telephone calls not captured.}  A significant portion of
          service interactions is conducted over the phone and is not recorded
          in the ticket system.  LLM extraction is therefore limited to what
          is present in the digital record and may miss important problem
          context that was communicated verbally.

    \item \textbf{Zero-shot only.}  No fine-tuning or few-shot examples were
          used in this study.  Fine-tuned smaller models could potentially
          match or exceed the performance of the larger zero-shot models at
          significantly lower inference cost.

    \item \textbf{Benchmark scope.}  The benchmark covers three ticket
          categories and 250~tickets.  Generalisation to other ticket types
          in the same system, or to service ticket systems in different
          industrial domains, is not guaranteed without further evaluation.

    \item \textbf{API cost and availability.}  Proprietary API models~(GPT-4o)
          introduce per-token costs and depend on external service availability,
          which limits their suitability for a high-volume production
          deployment where the Ollama-hosted models have a clear advantage.

    \item \textbf{Context-window utilisation.}  Very long email threads~(20+
          messages) may saturate even the 128k context window for the smallest
          models.  Truncation strategies were not systematically evaluated,
          and their impact on extraction quality is unknown.
\end{itemize}
